\chapter{Understanding heritage speakers’ languages: Conclusion and outlook}\label{chap:key:8}

This study investigated heritage speakers’ languages by examining various factors that contribute to differences between heritage and majority languages, with a particular focus on word order in Turkish and Kurmanji in Germany and Turkey. As outlined in the introduction, research on heritage languages lacks linguistic diversity, with the majority of existing studies concentrating on languages such as Spanish and Chinese in contact with English. By focusing on Turkish and Kurmanji, the present study contributes to broadening the empirical scope of heritage language research.

While Turkish as a heritage language has received considerable attention, particularly in the German context, heritage Kurmanji remains virtually unexplored. Research on Kurmanji-Turkish bilinguals is also extremely limited, both with respect to their Kurmanji and their Turkish. The present study addresses this gap by providing the first systematic investigation of Kurmanji as a heritage language in Germany and by examining both Kurmanji and Turkish among Kurmanji-Turkish bilingual speakers.

Although \citegen{Rothman2009} definition of a heritage language encompasses a broad spectrum of speakers and languages worldwide, research on heritage languages is often limited to immigration contexts, particularly in Europe and North America. By including Kurmanji-Turkish bilinguals, the study challenged the often implicit assumption that heritage speakers are primarily associated with immigrant languages in the Global North. It highlighted the importance of recognizing linguistic diversity among heritage speakers, including speakers of minoritized languages.

The study built on existing proposals that aim to explain why heritage languages differ from majority languages, with a particular focus on transfer and divergent attainment accounts. To explore the divergent attainment account more comprehensively, the study examined not only commonly investigated factors such as cross-linguistic influence and bilingualism, but also the societal status of the languages, internal linguistic dynamics, particularly concentrating on the role of register.

Heritage language research often compares bilingual heritage speakers with monolingual majority speakers, thereby investigating two factors simultaneously: bilingualism and the societal status of the language. To disentangle these factors, the study included an additional group of speakers—Kurmanji-Turkish bilinguals in Turkey. While these speakers are bilingual, like heritage speakers, Turkish in this context is not a heritage language but the majority language.

Adding this group introduced another parameter: the intensity of contact. In Germany, contact between Turkish or Kurmanji and the dominant German language began in the late 20th century with the arrival of ‘guest workers,’ while Turkish-Kurmanji contact dates back to the Ottoman Empire and intensified with the establishment of the Turkish Republic (Yağmur 2001). The intensity of contact appears to play a significant role in the factors driving differences between the groups. Transfer was observed only in the Kurmanji spoken in Turkey, where language contact has been ongoing for years. In contrast, no trace of cross-linguistic influence was found in heritage Turkish or heritage Kurmanji spoken in Germany.

That said, the statement that no transfer is observed in more recent contact situations refers, of course, only to the phenomenon investigated in this study and does not rule out the possibility of transfer from the dominant language in more recent contact situations. Undoubtedly, such transfer exists, and extensive research provides ample evidence of this phenomenon, as shown in Chapter \ref{chap:key:1}.

The key point here is that other factors must also be explored. This study specifically highlighted the role of internal dynamics in the language, including changes in the (monolingual) majority variety that are unrelated to language contact. To observe such dynamics, however, it is essential to collect and analyse the informal speech of the majority language speakers. For example, the study revealed that noncanonical phenomena, such as the placement of new referents in the post-predicate position, are also found in the informal registers of majority Turkish speakers. Without collecting data from different registers, it would not have been possible to attest these ongoing differences in the monolingual variety.

Second, incorporating register into the analysis of heritage language use made it possible to identify more precisely the contexts in which heritage speakers diverge from majority language speakers. When the frequency of placing new information in the post-predicate position was examined without taking register differences into account, the results suggested that this noncanonical phenomenon occurs much more frequently among heritage speakers than among majority speakers. However, a closer examination of new information placement across different communicative situations revealed that this apparent frequency difference stems from a lack of register flexibility in heritage speakers. Specifically, while majority speakers tend to place new information in the post-predicate position only in informal settings, heritage speakers do so across all communicative situations, without distinguishing between formal and informal settings. Such a finding is unsurprising, given the conditions under which heritage languages are acquired—primarily at home, with limited or no access to formal education in the language.

The finding that differences between heritage and majority speakers primarily manifest in formal registers supports the proposal centered on divergent attainment in heritage speakers. Heritage speakers acquire their language in a context different from that of majority speakers, often without formal education in the language. As a result, they exhibit less flexibility in formal registers. This finding also emphasizes the crucial need to provide opportunities for heritage speakers to engage with the formal registers of their heritage language and highlights the importance of ensuring that heritage languages are included in school curricula, expanding the number of language classes available, developing textbooks, and creating all possibilities for learning all languages, not just the ones that are traditionally considered prestigious. Implementing such measures would foster a more inclusive and equitable language education that fully embraces linguistic diversity.

Furthermore, the study also revealed that register inflexibility concerns not only heritage speakers but also speakers of the majority languages. However, unlike the extension of features typical of informal language into formal situations observed in heritage speakers, majority language speakers tend to employ more formal language even in informal situations. Such register inflexibility may stem from language anxiety experienced by bilingual speakers, even when communicating in their majority language (\citealt{SevincDewaele2018}, \citealt{Bunk2024}). The reason for this is strong monolingual habitus which takes monolingualism as the norm and does not welcome any foreign influence on the language (\citealt{Wiese2020}, \citealt{KerswillWiese2022}). The register inflexibility attested in the majority language emphasizes the importance of embracing multilingualism and abandoning the idea that monolingualism is the norm, especially considering that the majority of the world’s population is multilingual.

Future research on heritage speakers should consider register as one of the factors and explore other linguistic domains and language phenomena to determine whether what was previously attributed to cross-linguistic influence could, in fact, stem from register inflexibility in bilinguals. Much of the research conducted in the RUEG project has already addressed this (\citealt{AlexiadouRizou2023}, \citealt{WieseEtAl2022}, \citealt{Tsehaye2023}, \citealt{SchroederEtAl2024}). Additionally, the RUEG corpora (\citealt{LudelingEtAl2022}, \citealt{IefremenkoEtAl2024}) available online can be used for further research, providing an opportunity to investigate the role of register in bilinguals.

Furthermore, little is known about the role register plays in the processing of language by speakers. A limited but emerging body of research has examined how monolingual speakers process (mis)matches in formality/register (\citealt{MaquateEtAl2025}, \citealt{PescumaEtAl2024}), yet comparable evidence for bilingual populations remains largely absent. This gap is currently being addressed within the Collaborative Research Centre 1412 “Register: Language Users’ Knowledge of Situational-Functional Variation” (CRC 1412), a DFG-funded research centre based in Germany. The CRC investigates register from multiple linguistic perspectives; within this framework, one subproject focuses specifically on the real-time processing of register mismatches among heritage speakers in Germany.

At the same time, it is important to acknowledge that not every heritage language can be explored through such a broad range of factors. As mentioned earlier, a heritage language is not necessarily one spoken by the children or grandchildren of immigrants to the Global North. Many minoritized languages exhibit the characteristics of a heritage language. However, such languages may lack a monolingual variety because often speakers of these minority languages tend to be bilingual or even multilingual, with a dominant proficiency in the societal language of the region they live in. As a result, these minority languages typically do not have a majority variety, and consequently, may not possess a wide range of registers, or at least speakers are often not exposed to or do not have opportunities to engage with these different registers.

Recognizing these complexities highlights the importance of broadening the scope of heritage language research. To truly understand heritage speakers' languages and model their grammar, it is crucial to include a diverse range of languages and language pairs. This approach would foster dialogue between studies on minoritized languages and heritage language research, fields that, while addressing similar issues, often operate in isolation from one another. Building on these considerations, the present study took a step toward bridging this gap by positioning Kurmanji as a heritage language in Turkey and contributed to research on both heritage Kurmanji and majority Turkish among Kurmanji-Turkish bilinguals. Furthermore, by making the corpus data openly accessible, the study aimed to encourage further research on these language pairs and to foster broader engagement with languages that have traditionally been underrepresented in heritage language research.
